% Generated draft: requires mathematical and source review.
\documentclass[11pt]{article}
\usepackage[margin=25mm]{geometry}
\usepackage{fontspec}
\setmainfont{Latin Modern Roman}
\usepackage{amsmath,amssymb,mathtools,unicode-math}
\setmathfont{Latin Modern Math}
\usepackage{enumitem,needspace,xurl,hyperref}
\setlength{\parindent}{0pt}
\setlength{\parskip}{5pt}
\setlength{\emergencystretch}{3em}
\newcommand{\N}{\mathbb{N}}
\newcommand{\Z}{\mathbb{Z}}
\newcommand{\Q}{\mathbb{Q}}
\newcommand{\R}{\mathbb{R}}
\newcommand{\C}{\mathbb{C}}
\newcommand{\F}{\mathbb{F}}
\newcommand{\A}{\mathbb{A}}
\newcommand{\PP}{\mathbb{P}}
\newcommand{\E}{\mathbb{E}}
\newcommand{\eps}{\varepsilon}
\newcommand{\dd}{\,\mathrm{d}}
\begin{document}
\section*{UnsolvedMath ID: 30003570}
\textbf{Source problem code:} catalogue entry
\subsection{Definitions and mathematical statement}
Write \(X_d=\mathcal F_{2d}\) for the space intended in the source. The supplied record does not define this space, the admissible values of \(d\), or the meaning of ``primitive Heegner divisor.'' Consequently, the following is a conditional formulation; these missing specifications must be supplied before it becomes a fully determined mathematical problem.

Fix an admissible parameter \(d\), a common finite-dimensional real vector space \(V_d\) of divisor classes on \(X_d\), and a class map \(D\mapsto[D]\in V_d\). Let \(\mathscr H_d^{\mathrm{prim}}\) denote the intended family of effective primitive Heegner divisors. Define its generated cone by finite nonnegative real combinations:
\[
 C_d^{\mathrm{prim}}
 =\left\{\sum_{i=1}^{r}a_i[H_i]:
 r\in\Z_{\geq0},\ a_i\in\R_{\geq0},\ H_i\in\mathscr H_d^{\mathrm{prim}}\right\}.
\]
Using the same divisor-class convention, define the pseudo-effective cone by
\[
 \overline{\operatorname{Eff}}(X_d)
 =\overline{\operatorname{cone}\{[D]:D\text{ is an effective divisor on }X_d\}}^{\,V_d},
\]
where the bar denotes closure in the usual topology of \(V_d\). This closed cone is the object called \(\operatorname{Eff}(\mathcal F_{2d})\), the pseudo-effective cone, in the supplied question.

The requested comparison is
\[
 C_d^{\mathrm{prim}}\stackrel{?}{=}
 \overline{\operatorname{Eff}}(X_d),
 \qquad\text{or}\qquad
 C_d^{\mathrm{prim}}\subsetneq
 \overline{\operatorname{Eff}}(X_d).
\]
The inclusion presupposes that the specified Heegner divisors are effective and that both cones use the same class map. The source does not say whether the generated Heegner cone itself should be closed; replacing it by \(\overline{C_d^{\mathrm{prim}}}\) would give a potentially different question.

The notation suggests a moduli space of polarized K3 surfaces of degree \(2d\), but the supplied record does not specify whether polarizations are primitive, ample, or merely quasi-polarizations, nor whether a coarse space, stack, or compactification is intended. These choices, the divisor equivalence relation, and any boundary divisors must be fixed. The record also does not specify whether the comparison concerns a particular \(d\), every admissible \(d\), or dependence on \(d\). No choice among these interpretations is asserted here.
\subsection{Short English statement}
The question asks whether every pseudo-effective divisor class on the intended space \(\mathcal F_{2d}\) belongs to the cone generated by primitive Heegner divisor classes, or whether that cone is strictly smaller. The moduli space, degree range and quantifier, divisor-class convention, meaning of primitive Heegner divisor, and closure convention require clarification.
\subsection{Sources}
\begin{itemize}
\item Local registry: \texttt{data/unsolvedmath/problems.json} and supplementary \texttt{catalog/problems.json}; supplied record with permanent UnsolvedMath ID \(30003570\), exact source code \texttt{catalogue entry}, title ``Heegner Divisors and the Pseudo-Effective Cone.'' These are registry references supplied for identification, not independently inspected files.
\item Supplied report citation: \emph{Komplexe Analysis} (2018), \emph{Oberwolfach Reports}, pp.\ 2427--2473, DOI \texttt{10.4171/owr/2017/39}. \url{https://doi.org/10.4171/owr/2017/39}. The record attributes the problem to Oberwolfach workshop open-problem discussions, gives proposed year 2017, and names no proposer. The citation's 2018 year is retained alongside the 2017 DOI component and proposed year.
\item Supplied literature reference, labeled ``Bruinier--Moller'': \url{https://doi.org/10.1090/jag/734}.
\item Supplied literature reference, labeled ``Extremal divisors on K3 moduli'': \url{https://arxiv.org/abs/2504.16730}.
\item Supplied literature reference, labeled ``OWR report'': \url{https://doi.org/10.4171/owr/2017/39}.
\item Literature assessment dated exactly \texttt{2026-08-21}: the supplied curation labels the problem open and states that equality of the cone generated by effective Heegner divisors with the full pseudo-effective cone remains open, while finite generation and substantial arithmetic cases are known. This is supplied dated curation derived from the OWR record, not an independently verified literature assessment. Its claims and time-sensitive status require independent verification; its wording ``effective Heegner divisors'' also requires reconciliation with ``primitive Heegner divisors'' in the question.
\end{itemize}
\end{document}
